\section{AMC 8 1996}
        \begin{enumerate}
        	\item (\textit{AMC 8 1996}) How many positive factors of 36 are also multiples of 4?


 $\text{(A)}\ 2 \qquad \text{(B)}\ 3 \qquad \text{(C)}\ 4 \qquad \text{(D)}\ 5 \qquad \text{(E)}\ 6$

 \textbf{Answer: }B



\textbf{Solution}

The factors of $36$ are $1, 2, 3, 4, 6, 9, 12, 18,$ and $36$.

 The multiples of $4$ up to $36$ are $4, 8, 12, 16, 20, 24, 28, 32$ and $36$.

 Only $4, 12$ and $36$ appear on both lists, so the answer is $3$, which is option $\boxed{B}$.

 



\textbf{Solution 2}

$36 = 4^1 \cdot 3^2$. All possible factors of $36$ will be here, except for ones divisible by $2$ and not by $4$. $(1+1)\cdot (2+1) = 6$. Subtract factors not divisible by $4$, which are $1$, $3^1$, and $3^2$. $6-3=3$, which is $\boxed{B}$.

 



\textbf{Solution 3}

Divide $36$ by $4$, and the remaining factors, when multiplied by $4$, will be factors of $36$.

 $36 \div 4 = 9$, which has $3$ factors, giving us option $\boxed{B}$.

 


	\item (\textit{AMC 8 1996}) Jose, Thuy, and Kareem each start with the number 10.  Jose subtracts 1 from the number 10, doubles his answer, and then adds 2.  Thuy doubles the number 10, subtracts 1 from her answer, and then adds 2.  Kareem subtracts 1 from the number 10, adds 2 to his number, and then doubles the result.  Who gets the largest final answer?


 $\text{(A)}\ \text{Jose} \qquad \text{(B)}\ \text{Thuy} \qquad \text{(C)}\ \text{Kareem} \qquad \text{(D)}\ \text{Jose and Thuy} \qquad \text{(E)}\ \text{Thuy and Kareem}$

 \textbf{Answer: }C



\textbf{Solution}

Jose gets $10 - 1 = 9$, then $9 \cdot 2 = 18$, then $18 + 2 = 20$.

 Thuy gets $10 \cdot 2 = 20$, then $20 - 1 = 19$, and then $19 + 2 = 21$.

 Kareem gets $10 - 1 = 9$, then $9 + 2 = 11$, and then $11\cdot 2 = 22$.

 Thus, Kareem gets the highest number, and the answer is $\boxed{C}$.

 


 


	\item (\textit{AMC 8 1996}) The 64 whole numbers from 1 through 64 are written, one per square, on a checkerboard (an 8 by 8 array of 64 squares).  The first 8 numbers are written in order across the first row, the next 8 across the second row, and so on.  After all 64 numbers are written, the sum of the numbers in the four corners will be


 $\text{(A)}\ 130 \qquad \text{(B)}\ 131 \qquad \text{(C)}\ 132 \qquad \text{(D)}\ 133 \qquad \text{(E)}\ 134$

 \textbf{Answer: }A



\textbf{Solution}

Obviously $1$ is in the top left corner, $8$ is in the top right corner, and $64$ is in the bottom right corner. To find the bottom left corner, subtract $7$ from $64$ which is $57$. Adding the results gives $1+8+57+64=130$ which is answer $\boxed{A}$.

 


	\item (\textit{AMC 8 1996}) $\dfrac{2+4+6+\cdots + 34}{3+6+9+\cdots+51}=$
$\text{(A)}\ \dfrac{1}{3} \qquad \text{(B)}\ \dfrac{2}{3} \qquad \text{(C)}\ \dfrac{3}{2} \qquad \text{(D)}\ \dfrac{17}{3} \qquad \text{(E)}\ \dfrac{34}{3}$

 \textbf{Answer: }B



\textbf{Solution 1}

First, notice that each number in the numerator is a multiple of $2$, and each number in the denominator is a multiple of $3$.  This suggests that each expression can be factored.  Factoring gives:

 $\dfrac{2(1+2+3+\cdots + 17)}{3(1+2+3+\cdots + 17)}$

 Since all the material in the parentheses is the same, the common factor in the numerator and the denominator may be cancelled, leaving $\frac{2}{3}$, which is option $\boxed{B}$.

 



\textbf{Solution 2}

There are $17$ terms in the numerator $2 + 4 + 6 + ... + 30 + 32+ 34$.  Consider adding those terms, but in a different order.  Start with the last two terms, $2 + 34$, and then add the next two terms on the outside, $4 + 32$, and continue.  You will get $8$ pairs of numbers that add to $36$, while the $9^{th}$ number in the middle will be alone.  That number is $18$.  Adding all the numbers gives $8\cdot 36 + 18 = 306$.

 Similarly, the denominator has $17$ terms of $3 + 6 + 9 + ... + 45 + 48 + 51$.  There are $8$ pairs of numbers that add up to $54$, with the $9^{th}$ number in the center being $27$.  The total of all the numbers is $8\cdot 54 + 27 = 459$.

 The answer is $\frac{306}{459}$.  Eyeballing the options, the fraction is clearly under $1$, but more than $\frac{1}{2}$.  Thus, the answer must be $\frac{2}{3}$, or $\boxed{B}$.  Alternately, you can do the work by factoring out a $9$ in the numerator to give $\frac{34}{51}$.  Factoring out $17$ will give the desired answer.

 



\textbf{Solution 3}

Start by finding a pattern:

 $\frac{2}{3} = \frac{2}{3}$
$\frac{2 + 4}{3 + 6} = \frac{6}{9} = \frac{2}{3}$
$\frac{2 + 4 + 6}{3 + 6 + 9} = \frac{12}{18} = \frac{2}{3}$

 Each step doesn't seem to change the value of the fraction, so $\boxed{B}$ is the right answer.

 


	\item (\textit{AMC 8 1996}) The letters $P$, $Q$, $R$, $S$, and $T$ represent numbers located on the number line as shown.


 
                    \begin{figure}[H]
                        \centering
                        \includegraphics[width=0.5\textwidth]{Problem/AMC_8_AJHSME_1996_5_Problem_Figure_1.png}
                    \end{figure}
                    


 Which of the following expressions represents a negative number?


 $\text{(A)}\ P-Q \qquad \text{(B)}\ P\cdot Q \qquad \text{(C)}\ \dfrac{S}{Q}\cdot P \qquad \text{(D)}\ \dfrac{R}{P\cdot Q} \qquad \text{(E)}\ \dfrac{S+T}{R}$

 \textbf{Answer: }A



\textbf{Solution 1}

First, note that $P < Q < 0 < R < S < T$.  Thus, $P$ and $Q$ are negative, while $R$, $S$, and $T$ are positive.

 Option $A$ is the difference of two numbers.  If the first number is lower than the second number, the answer will be negative.  In this case, $P < Q$, so $P - Q < 0$, and $P-Q$ is thus indeed a negative number.  So option $\boxed{A}$ is correct.

 Option $B$ is the product of two negatives, which is always positive.

 Option $C$ starts with the quotient of a positive and a negative.  This is negative.  But this negative quotient is then multiplied by a negative number, which gives a positive number.

 Option $D$ has the product of two negatives in the denominator.  Thus, the denominator is positive.  Additionally, the numerator is (barely) positive.  A positive over a positive will be positive.

 Option $E$ has the sum of two positives in the numerator.   That will be positive.  The denominator is also positive.  Again, the quotient of two positives is positive.

 



\textbf{Solution 2}

Estimate the numbers, and calculate.

 $P\approx -3.5$
$Q\approx -1.1$
$R\approx 0.1$
$S\approx 0.9$
$T\approx 1.5$

 For option $A$, the value of $P - Q = -3.5 - (-1.1) = -3.5 + 1.1 = -2.4$.  Thus $\boxed{A}$ is the right answer.

 For option $B$, $PQ = -3.5 \cdot -1.1 = +3.85$

 For option $C$, $\frac{S}{Q}\cdot P = \frac{0.9}{-1.1}\cdot -3.5 \approx -0.818 \cdot -3.5 \approx +2.836$

 For option $D$, $\frac{R}{PQ} = \frac{0.1}{3.85} \approx +0.026$

 For option $E$, $\frac{S+T}{R} = \frac{0.9 + 1.5}{0.1} = \frac{2.4}{0.1} = +24$



	\item (\textit{AMC 8 1996}) What is the smallest result that can be obtained from the following process?


 
            \begin{itemize}
        \item Choose three different numbers from the set $\{3,5,7,11,13,17\}$.

        \item Add two of these numbers.

        \item Multiply their sum by the third number.
            \end{itemize}
            $\text{(A)}\ 15 \qquad \text{(B)}\ 30 \qquad \text{(C)}\ 36 \qquad \text{(D)}\ 50 \qquad \text{(E)}\ 56$

 \textbf{Answer: }C



\textbf{Solution}

Since we want the smallest possible result, and we are only adding and multiplying positive numbers over $1$, we can "prune" the set to the three smallest numbers $\{3,5,7\}$.  Using bigger numbers will create bigger sums and bigger products.

 From there, compute the $3$ ways you can do the two operations:

 $(3+5)7 = 8\cdot 7 = 56$
$(3 + 7)5 = 10\cdot 5 = 50$
$(7+5)3 = 12\cdot 3 = 36$

 The smallest number is 36, giving an answer of $\boxed{C}$



	\item (\textit{AMC 8 1996}) Brent has goldfish that quadruple (become four times as many) every month, and Gretel has goldfish that double every month.  If Brent has 4 goldfish at the same time that Gretel has 128 goldfish, then in how many months from that time will they have the same number of goldfish?


 $\text{(A)}\ 4 \qquad \text{(B)}\ 5 \qquad \text{(C)}\ 6 \qquad \text{(D)}\ 7 \qquad \text{(E)}\ 8$

 \textbf{Answer: }B



\textbf{Solution 1}

Call this month "Month 0".  Make a table of the fish that Brent and Gretel have each month.

 $\text{Month / Brent / Gretel}$
$\text{0 / 4 / 128}$
$\text{1 / 16 / 256}$
$\text{2 / 64 / 512}$
$\text{3 / 256 / 1024}$
$\text{4 / 1024 / 2048}$
$\text{5 / 4096 / 4096}$

 You could create a similar table without doing all of the calculations by converting all the goldfish into powers of $2$.  In this table, you could increase Brent's goldfish by two powers of $2$, while increasing Gretel's fish by one power of $2$.

 Either way, in $5$ months they will have the same number of fish, giving an answer of $\boxed{B}$




\textbf{Solution 2}

Create two equations.

 Brent starts with $4$ goldfish and they quadruple each month, giving us $y=4\cdot 4^{x}$ or $y=2^{2x+2}$.

 Gretel starts with $128$ goldfish and they double each month, giving us $y=128 \cdot 2^{x}$ or $y=2^{x+7}$.

 Using substitution, $2^{2x+2}=2^{x+7}$. Thus $2x+2=x+7$ giving us $x=5$, which is $\boxed{B}$.

 


	\item (\textit{AMC 8 1996}) Points $A$ and $B$ are 10 units apart.  Points $B$ and $C$ are 4 units apart.  Points $C$ and $D$ are 3 units apart.  If $A$ and $D$ are as close as possible, then the number of units between them is


 $\text{(A)}\ 0 \qquad \text{(B)}\ 3 \qquad \text{(C)}\ 9 \qquad \text{(D)}\ 11 \qquad \text{(E)}\ 17$

 \textbf{Answer: }B



\textbf{Solution}

If $AB = 10$ and $BC=4$, then $(10 - 4) \le AC \le (10 + 4)$ by the \href{https://artofproblemsolving.com/wiki/index.php?title=Triangle\_inequality}{triangle inequality}.  In the triangle inequality, the equality is only reached when the "triangle" $ABC$ is really a degenerate triangle, and $ABC$ are collinear.  

 Simplifying, this means the smallest value $AC$ can be is $6$.  

 Applying the triangle inequality on $ACD$ with $AC = 6$ and $CD = 3$, we know that $6 - 3 \le AD \le 6 + 3$ when $AC$ is minimized.  If $AC$ were larger, then $AD$ could be larger, but we want the smallest $AD$ possible, and not the largest.  Thus, $AD$ must be at least $3$, but cannot be smaller than $3$.  Therefore, $\boxed{B}$ is the answer.

 This answer comes when $ABCD$ are all on a line, with $A=0, B=10, C = 6,$ and $D=3$.

 


	\item (\textit{AMC 8 1996}) If 5 times a number is 2, then 100 times the reciprocal of the number is


 $\text{(A)}\ 2.5 \qquad \text{(B)}\ 40 \qquad \text{(C)}\ 50 \qquad \text{(D)}\ 250 \qquad \text{(E)}\ 500$

 \textbf{Answer: }D



\textbf{Solution}

If $5$ times a number is $2$, then $5x = 2 \rightarrow x = \frac{2}{5}$, and the number is $\frac{2}{5}$.

 If the number is $\frac{2}{5}$, then $100$ times its reciprocal is $100$ times $\frac{5}{2}$, which is $250$, giving an answer of $\boxed{D}$.

 


	\item (\textit{AMC 8 1996}) When Walter drove up to the gasoline pump, he noticed that his gasoline tank was 1/8 full.  He purchased 7.5 gallons of gasoline for $\$10$.  With this additional gasoline, his gasoline tank was then 5/8 full.  The number of gallons of gasoline his tank holds when it is full is


 $\text{(A)}\ 8.75 \qquad \text{(B)}\ 10 \qquad \text{(C)}\ 11.5 \qquad \text{(D)}\ 15 \qquad \text{(E)}\ 22.5$

 \textbf{Answer: }D



\textbf{Solution}

The tank started at $\frac{1}{8}$ full, and ended at $\frac{5}{8}$ full.  Therefore, Walter filled $\frac{5}{8} - \frac{1}{8} = \frac{4}{8} = \frac{1}{2}$ of the tank.

 If Walter fills half the tank with $7.5$ gallons, then Walter can fill two halves of the tank (or a whole tank) with $7.5 \times 2 = 15$ gallons, giving an answer of $\boxed{D}$



	\item (\textit{AMC 8 1996}) Let $x$ be the number

 \[0. \underbrace{0000...0000}_{1996\text{ zeros}}1,\]
where there are 1996 zeros after the decimal point.  Which of the following expressions represents the largest number?


 $\text{(A)}\ 3+x \qquad \text{(B)}\ 3-x \qquad \text{(C)}\ 3\cdot x \qquad \text{(D)}\ 3/x \qquad \text{(E)}\ x/3$

 \textbf{Answer: }D



\textbf{Solution}

Estimate each of the options.

 $A$ will give a number that is just over $3$.

 $B$ will give a number that is just under $3$.  This eliminates $B$, because $A$ is bigger.

 $C$ will give a number that is barely over $0$, since it is three times a tiny number.  This eliminates $C$, because $A$ is bigger.

 $D$ will give a huge number.  $\frac{1}{x}$ will get very, very large in magnitude when $x$ gets close to zero.  You can see this by examining the sequence $x=0.1$, which gives $10$ as the reciprocal, $x=0.01$, which gives $100$ as the reciprocal, and $x=0.001$, which gives $1000$ as the reciprocal.  Thus, $D$ will be huge, and this eliminates $A$.

 $E$ will give a small number, since you're dividing a tiny number into thirds.  This eliminates $E$, and gives $\boxed{D}$ as the answer.

 


	\item (\textit{AMC 8 1996}) What number should be removed from the list

 \[1,2,3,4,5,6,7,8,9,10,11\]
so that the average of the remaining numbers is $6.1$?


 $\text{(A)}\ 4 \qquad \text{(B)}\ 5 \qquad \text{(C)}\ 6 \qquad \text{(D)}\ 7 \qquad \text{(E)}\ 8$

 \textbf{Answer: }B



\textbf{Solution 1}

Adding all of the numbers gives us $\frac{11\cdot12}{2}=66$ as the current total. Since there are $11$ numbers, the current average is $\frac{66}{11}=6$. We need to take away a number from the total and then divide the result by $10$ because there will only be $10$ numbers left to give an average of $6.1$. Setting up the equation:

 $\frac{66-x}{10}=6.1$
$66 - x = 61$
$x = 5$

 Thus, the answer is $\boxed{B}$




\textbf{Solution 2}

Similar to the first solution, the current total is $66$.  Since there are $11$ numbers on the list, taking $1$ number away will leave $10$ numbers.  If those $10$ numbers have an average of $6.1$, then those $10$ numbers must have a sum of $10 \times 6.1 = 61$.  Thus, the number that was removed must be $66 - 61 = 5$, and the answer is $\boxed{B}$.

 


	\item (\textit{AMC 8 1996}) In the fall of 1996, a total of 800 students participated in an annual school clean-up day.  The organizers of the event expect that in each of the years 1997, 1998, and 1999, participation will increase by 50\% over the previous year.  The number of participants the organizers will expect in the fall of 1999 is


 $\text{(A)}\ 1200 \qquad \text{(B)}\ 1500 \qquad \text{(C)}\ 2000 \qquad \text{(D)}\ 2400 \qquad \text{(E)}\ 2700$

 \textbf{Answer: }E



\textbf{Solution 1}

If the participation increases by $50\%$, then it is the same as saying participation is multipled by a factor of $100\% + 50\% = 1 + 0.5 = 1.5$.

 In 1997, participation will be $800 \cdot 1.5 = 1200$.

 In 1998, participation will be $1200 \cdot 1.5 = 1800$

 In 1999, participation will be $1800 \cdot 1.5 = 2700$, giving an answer of $\boxed{E}$.

 



\textbf{Solution 2}

Since the percentage increase is the same each year, this is an example of exponential growth with a base of $1.5$.  In three years, there will be $1.5^3 = \frac{27}{8}$ times as many participants.  Multiplying this by the $800$ current participants, there are $2700$ participants, and the answer is $\boxed{E}$.

 


	\item (\textit{AMC 8 1996}) Six different digits from the set

 \[\{ 1,2,3,4,5,6,7,8,9\}\]
are placed in the squares in the figure shown so that the sum of the entries in the vertical column is 23 and the sum of the entries in the horizontal row is 12.
The sum of the six digits used is


 
                    \begin{figure}[H]
                        \centering
                        \includegraphics[width=0.5\textwidth]{Problem/AMC_8_AJHSME_1996_14_Problem_Figure_1.png}
                    \end{figure}
                    
$\text{(A)}\ 27 \qquad \text{(B)}\ 29 \qquad \text{(C)}\ 31 \qquad \text{(D)}\ 33 \qquad \text{(E)}\ 35$

 \textbf{Answer: }B



\textbf{Solution}

Looking at the vertical column, the three numbers sum to $23$.  If we make the numbers on either end $9$ and $8$ in some order, the middle number will be $6$.  This is the minimum for the intersection.

 Looking at the horizontal row, the four numbers sum to $12$.  If we minimize the three numbers on the right to $123$, the first number has a maximum value of $6$.  This is the maximum for the intersection

 Thus, the minimum of the intersection is $6$, and the maximum of the intersection is $6$.  This means the intersection must be $6$, and the other numbers must be $9$ and $8$ in the column, and $123$ in the row.  The sum of all the numbers is $12 + 23 - 6 = \boxed{\text{(B)}\ 29}$.

 


 


	\item (\textit{AMC 8 1996}) The remainder when the product $1492\cdot 1776\cdot 1812\cdot 1996$ is divided by 5 is


 $\text{(A)}\ 0 \qquad \text{(B)}\ 1 \qquad \text{(C)}\ 2 \qquad \text{(D)}\ 3 \qquad \text{(E)}\ 4$

 \textbf{Answer: }E



\textbf{Solution 1}

To determine a remainder when a number is divided by $5$, you only need to look at the last digit.  If the last digit is $0$ or $5$, the remainder is $0$.  If the last digit is $1$ or $6$, the remainder is $1$, and so on.

 To determine the last digit of $1492\cdot 1776\cdot 1812\cdot 1996$, you only need to look at the last digit of each number in the product.  Thus, we compute $2\cdot 6\cdot 2\cdot 6 = 12^2 = 144$.  The last digit of the number $1492\cdot 1776\cdot 1812\cdot 1996$ is also $4$, and thus the remainder when the number is divided by $5$ is also $4$, which gives an answer of $\boxed{E}$.

 



\textbf{Solution 2}

$1492\cdot 1776\cdot 1812\cdot 1996 \mod 5$
$\equiv 2\cdot 1\cdot 2\cdot 1 \mod 5$
$=4$, which is option $\boxed{E}$.

 


	\item (\textit{AMC 8 1996}) $1-2-3+4+5-6-7+8+9-10-11+\cdots + 1992+1993-1994-1995+1996=$
$\text{(A)}\ -998 \qquad \text{(B)}\ -1 \qquad \text{(C)}\ 0 \qquad \text{(D)}\ 1 \qquad \text{(E)}\ 998$

 \textbf{Answer: }C



\textbf{Solution}

Put the numbers in groups of $4$:

 $(1-2-3+4)+(5-6-7+8)+(9-10-11+ 12) + \cdots + (1993-1994-1995+1996)$

 The first group has a sum of $0$.

 The second group increases the two positive numbers on the end by $1$, and decreases the two negative numbers in the middle by $1$.  Thus, the second group also has a sum of $0$.

 Continuing the pattern, every group has a sum of $0$, and thus the entire sum is $0$, giving an answer of $\boxed{C}$.

 



\textbf{Solution 2}

Let any term of the series be $t_n$. Realize that at every $n\equiv0 \pmod4$, the sum of the series is 0. For $t_{1996}$ we know $1996\equiv0 \pmod4$ so the solution is $\boxed{C}$.

 



\textbf{Solution 3}

This solution is slightly more detailed than Solution 2, even though they are essentially the same.

 Compute the sum of the first and last term, the second and second last term, the third and third last term. We get $(1+1996)+(-2-1995)+(-3-1994)+(4+1993)$ and so on. We see that each pair's sum alternates between $1997$ and $-1997$.We see that every two consecutive pair cancels out. For example $(1+1996)+(-2-1995)$ gives a result of $0$. There are a total of $1996$ terms, there are $998$ pairs, and there are $499$ pairs which result in a sum of $1997$ each and $499$ pairs which result in a sum of $-1997$ each. These all cancel out to a sum of $0$, and thus the entire sum is $0$, giving an answer of $\boxed{C}$.

 


	\item (\textit{AMC 8 1996}) Figure $OPQR$ is a square.  Point $O$ is the origin, and point $Q$ has coordinates (2,2).  What are the coordinates for $T$ so that the area of triangle $PQT$ equals the area of square $OPQR$?


 
                    \begin{figure}[H]
                        \centering
                        \includegraphics[width=0.5\textwidth]{Problem/AMC_8_AJHSME_1996_17_Problem_Figure_1.png}
                    \end{figure}
                    


 NOT TO SCALE

 $\text{(A)}\ (-6,0) \qquad \text{(B)}\ (-4,0) \qquad \text{(C)}\ (-2,0) \qquad \text{(D)}\ (2,0) \qquad \text{(E)}\ (4,0)$

 \textbf{Answer: }C



\textbf{Solution}

The area of $\square OPQR$ is $2^2 = 4$.

 The area of $\triangle PQT$ is $\frac{1}{2}bh = \frac{1}{2}\cdot PT\cdot PQ$.

 If we set the areas equal, the area of $\triangle PQT$ is $4$.  Also, note that $PQ=2$.  Plugging those in, we get:

 $4 = \frac{1}{2} \cdot PT \cdot 2$
$PT = 4$

 If $PT = 4$, and $PO = 2$, then $OT = 2$, and $T$ must be $2$ units to the left of the origin.  This would be $(-2,0)$, giving answer $\boxed{C}$.

 


	\item (\textit{AMC 8 1996}) Ana's monthly salary was $2000 in May.  In June she received a 20% raise.  In July she received a 20% pay cut.  After the two changes in June and July, Ana's monthly salary was


 $\text{(A)}\ 1920\text{ dollars} \qquad \text{(B)}\ 1980\text{ dollars} \qquad \text{(C)}\ 2000\text{ dollars} \qquad \text{(D)}\ 2020\text{ dollars} \qquad \text{(E)}\ 2040\text{ dollars}$

 \textbf{Answer: }A



\textbf{Solution}

Notice that a 20\% raise translates to Ana's monthly salary multiplied by $(100 + 20)\% = 120\%$, and a 20\% pay cut translates to her salary multiplied by $100\% - 20\% = 80\%$, so

 

 \[2000 \cdot 120\% = 2400\]
 \[2400 \cdot 80\% = 1920 \rightarrow \boxed{\textbf{A}}\]



	\item (\textit{AMC 8 1996}) The pie charts below indicate the percent of students who prefer golf, bowling, or tennis at East Junior High School and West Middle School.  The total number of students at East is 2000 and at West, 2500.  In the two schools combined, the percent of students who prefer tennis is


 
                    \begin{figure}[H]
                        \centering
                        \includegraphics[width=0.5\textwidth]{Problem/AMC_8_AJHSME_1996_19_Problem_Figure_1.png}
                    \end{figure}
                    
$\text{(A)}\ 30\% \qquad \text{(B)}\ 31\% \qquad \text{(C)}\ 32\% \qquad \text{(D)}\ 33\% \qquad \text{(E)}\ 34\%$

 \textbf{Answer: }C



\textbf{Solution}

In the first school, $2000 \cdot 22\% = 2000 \cdot 0.22 = 440$ students prefer tennis.

 In the second school, $2500 \cdot 40\% = 2500 \cdot 0.40 = 1000$ students prefer tennis.

 In total, $440 + 1000 = 1440$ students prefer tennis out of a total of $2000 + 2500 = 4500$ students

 This means $\frac{1440}{4500}\cdot 100\% = \frac{32}{100} \cdot 100\% = 32\%$ of the students in both schools prefer tennis, giving answer $\boxed{\text{(C)}\ 32\%}$.

 


	\item (\textit{AMC 8 1996}) Suppose there is a special key on a calculator that replaces the number $x$ currently displayed with the number given by the formula $1/(1-x)$.  For example, if the calculator is displaying 2 and the special key is pressed, then the calculator will display -1 since $1/(1-2)=-1$.  Now suppose that the calculator is displaying 5.  After the special key is pressed 100 times in a row, the calculator will display 


 $\text{(A)}\ -0.25 \qquad \text{(B)}\ 0 \qquad \text{(C)}\ 0.8 \qquad \text{(D)}\ 1.25 \qquad \text{(E)}\ 5$

 \textbf{Answer: }A



\textbf{Solution}

We look for a pattern, hoping this sequence either settles down to one number, or that it forms a cycle that repeats.

 After $1$ press, the calculator displays $\frac{1}{1 - 5} = -\frac{1}{4}$

 After $2$ presses, the calculator displays $\frac{1}{1 - (-\frac{1}{4})} = \frac{1}{\frac{5}{4}} = \frac{4}{5}$

 After $3$ presses, the calculator displays $\frac{1}{1 - \frac{4}{5}} = \frac{1}{\frac{1}{5}} = 5$

 Thus, every three presses, the display will be $5$.  On press $3\cdot 33 = 99$, the display will be $5$.  One more press will give $-\frac{1}{4}$, which is answer $\boxed{A}$.

 


	\item (\textit{AMC 8 1996}) How many subsets containing three different numbers can be selected from the set 

 \[\{ 89,95,99,132, 166,173 \}\]
so that the sum of the three numbers is even?


 $\text{(A)}\ 6 \qquad \text{(B)}\ 8 \qquad \text{(C)}\ 10 \qquad \text{(D)}\ 12 \qquad \text{(E)}\ 24$

 \textbf{Answer: }D



\textbf{Solution}

To have an even sum with three numbers, we must add either $E+O+O$, or $E + E + E$, where $O$ represents an odd number, and $E$ represents an even number.

 Since there are not three even numbers in the given set, $E+E+E$ is impossible.  Thus, we must choose two odd numbers, and one even number.

 There are $2$ choices for the even number.

 There are $4$ choices for the first odd number.  There are $3$ choices for the last odd number.  But the order in which we pick these 2 numbers doesn't matter, so this overcounts the pairs of odd numbers by a factor of $2$.   Thus, we have $\frac{4\cdot 3}{2} = 6$ choices for a pair of odd numbers.

 In total, there are $2$ choices for an even number, and $6$ choices for the odd numbers, giving a total of $2\cdot 6 = 12$ possible choices for a 3-element set that has an even sum.  This is option $\boxed{D}$.

 


	\item (\textit{AMC 8 1996}) The horizontal and vertical distances between adjacent points equal 1 unit.  The area of triangle $ABC$ is


 
                    \begin{figure}[H]
                        \centering
                        \includegraphics[width=0.5\textwidth]{Problem/AMC_8_AJHSME_1996_22_Problem_Figure_1.png}
                    \end{figure}
                    
$\text{(A)}\ 1/4 \qquad \text{(B)}\ 1/2 \qquad \text{(C)}\ 3/4 \qquad \text{(D)}\ 1 \qquad \text{(E)}\ 5/4$

 \textbf{Answer: }B



\textbf{Solution 1}


                    \begin{figure}[H]
                        \centering
                        \includegraphics[width=0.5\textwidth]{Solution/AMC_8_AJHSME_1996_22_Solution_Figure_1.png}
                    \end{figure}
                    
$\triangle ADC$ takes up half of the 4x3 grid, so it has area of $6$.

 $\triangle ABD$ has height of $BF = 2$ and a base of $AD = 4$, for an area of $\frac{1}{2}\cdot 2 \cdot 4 = 4$.

 $\triangle CBD$ has height of $BE = 1$ and a base of $CD = 3$, for an area of $\frac{1}{2}\cdot 1 \cdot 3 = \frac{3}{2}$

 Note that $\triangle ABC$ can be found by taking $\triangle ADC$, and subtracting off $\triangle ABD$ and $\triangle CBD$.

 Thus, the area of $\triangle ABC = 6 - 4 - \frac{3}{2} = \frac{1}{2}$, and the answer is $\boxed{B}$.

 There are other equivalent ways of dissecting the figure; right triangles $\triangle ABF, \triangle BCE$ and rectangle $\square BEDF$ can also be used. You can also use $\triangle{BEC}$ and trapezoid $ADBE$.

 



\textbf{Solution 2}

Using the \href{https://artofproblemsolving.com/wiki/index.php?title=Shoelace\_Theorem}{Shoelace Theorem}, and labelling the points $A(0,0), B(3,2), C(4,3)$, we find the area is:

 $(0,0)$
$(3,2)$
$(4,3)$
$(0,0)$

 Area = $\frac{1}{2} = |(3\cdot 0 + 4\cdot 2 + 0\cdot 3) - (0\cdot 2 + 3\cdot 3 + 4\cdot 0)| = \frac{1}{2}$, which is option $\boxed{B}$.

 



\textbf{Solution 3}

Using \href{https://artofproblemsolving.com/wiki/index.php?title=Pick\%27s\_Theorem}{Pick's Theorem}, we can simply use the number of coordinate points to get the formula. In application, we get the formula $\frac{3}{2} + 0 - 1$, which equals $\frac{1}{2}$, giving us our answer $\boxed{B}$.

 



\textbf{Solution 4}

Using the Pythagorean Theorem, we know $BC$ is $\sqrt{2}$. Using $BC$ as a base of the triangle, the height is $\frac{1}{2}\sqrt{2}$. The area would be $\frac{1}{2}\cdot\frac{1}{2}\sqrt{2}\cdot\sqrt{2}$, which equals $\frac{1}{2}$, so the answer is $\boxed{B}$.

 


	\item (\textit{AMC 8 1996}) The manager of a company planned to distribute a $\$50$ bonus to each employee from the company fund, but the fund contained $\$5$ less than what was needed.  Instead the manager gave each employee a $\$45$ bonus and kept the remaining $\$95$ in the company fund.  The amount of money in the company fund before any bonuses were paid was


 $\text{(A)}\ \$945 \qquad \text{(B)}\ \$950 \qquad \text{(C)}\ \$955 \qquad \text{(D)}\ \$990 \qquad \text{(E)}\ \$995$

 \textbf{Answer: }E



\textbf{Solution 1}

Let $x$ be the number of employees at the company.  Then the total amount of money in the company fund, in dollars, can be written as either $50x-5$ or $45x+95$.  

 So $50x - 5 = 45x + 95$, and thus $x = 20$.  Substituting $x = 20$ into either $50x-5$ or $45x+95$ shows that the total amount of money in the company fund is $\boxed{\textbf{(E)}\ \$995}$.

 



\textbf{Solution 2 (Answer Choices)}

Since the number of employees is an integer, the amount of money in the fund, in dollars, must be $5$ less than a multiple of $50$.  This eliminates (B), (C), and (D).

 Also, the amount of money in the fund, in dollars, must also be $95$ more than a multiple of $45$.  This eliminates (A).

 Therefore, the correct answer is $\boxed{\textbf{(E)}\ \$995}$.

 


	\item (\textit{AMC 8 1996}) The measure of angle $ABC$ is $50^\circ$, $\overline{AD}$ bisects angle $BAC$, and $\overline{DC}$ bisects angle $BCA$.  The measure of angle $ADC$ is


 
                    \begin{figure}[H]
                        \centering
                        \includegraphics[width=0.5\textwidth]{Problem/AMC_8_AJHSME_1996_24_Problem_Figure_1.png}
                    \end{figure}
                    
$\text{(A)}\ 90^\circ \qquad \text{(B)}\ 100^\circ \qquad \text{(C)}\ 115^\circ \qquad \text{(D)}\ 122.5^\circ \qquad \text{(E)}\ 125^\circ$

 \textbf{Answer: }C



\textbf{Solution}

Let $\angle CAD = \angle BAD = x$, and let $\angle ACD = \angle BCD = y$

 From $\triangle ABC$, we know that $50 + 2x + 2y = 180$, leading to $x + y = 65$.

 From $\triangle ADC$, we know that $x + y + \angle D = 180$.  Plugging in $x + y = 65$, we get $\angle D = 180 - 65 = 115$, which is answer $\boxed{C}$.

 



\textbf{Solution 2}

Contruct $\overline{BE}$ through $D$ and intersects $\overline{AC}$ at point $E$

 By Exterior Angle Theorem,

 $\angle{ADE}$ $=$ $\angle{ABD} + \angle{BAD}$

 Similarly,

 $\angle{EDC} = \angle{DBC} + \angle{BCD}$

 Thus,

 $\angle{ADC} = \angle{ABC} + \angle{BAD} + \angle{BCD}$

 Let $\angle{ADC} = x$

 Because $\overline{AD}$ and $\overline{CD}$ are angle bisectors,

 
 \begin{align*} 180^\circ - x &= \angle{BAD} + \angle{BCD}\\ &=x - 50^\circ \end{align*}

$x = 115^\circ$

 Thus, the answer is $\boxed{C}$



	\item (\textit{AMC 8 1996}) A point is chosen at random from within a circular region.  What is the probability that the point is closer to the center of the region than it is to the boundary of the region?


 $\text{(A)}\ 1/4 \qquad \text{(B)}\ 1/3 \qquad \text{(C)}\ 1/2 \qquad \text{(D)}\ 2/3 \qquad \text{(E)}\ 3/4$

 \textbf{Answer: }A



\textbf{Solution 1}

Draw a circle with a radius of $2$.  Draw a concentric circle with radius $1$.  The circumference of this inner circle is the set of all points that are $1$ from the center, and $1$ from the outer circle meaning the set of points are equidistant from the center of the circles to the outside of the big circle.  

 The inside of smaller circle is the set of all points that are closer to the center of the region than to the boundary of the outer circle.  The "washer" region that is outside the circle of radius $1$, but inside the circle of radius $2$, is the set of all points that are closer to the boundary than to the center of the circle.

 If you select a random point in a region of area $B$, the probability that the point is in a smaller subregion $A$ is the ratio $\frac{A}{B}$.  In this case, $B = \pi\cdot 2^2 = 4\pi$, and $A = \pi\cdot 1^2 = \pi$, and the ratio of areas is $\frac{\pi}{4\pi} = \frac{1}{4}$, and the answer is $\boxed{A}$.

 



\textbf{Solution 1.1}

Alternatively, we could also let the radius of the entire circular region be $R$, the radius of the region where it is closer to the center be $r$.

 The probability is $\frac{1}{4}$, or answer choice $\boxed{A}$



\end{enumerate}